\documentclass[11pt,a4paper]{article}

% ---------------------------------------------------------------------------
% Encodage / langue
% ---------------------------------------------------------------------------
\usepackage[T1]{fontenc}
\usepackage[utf8]{inputenc}
\usepackage[french]{babel}
\usepackage{lmodern}

% ---------------------------------------------------------------------------
% Mise en page
% ---------------------------------------------------------------------------
\usepackage[margin=2.3cm]{geometry}
\usepackage{tabularx}
\usepackage{longtable}
\usepackage{booktabs}
\usepackage{array}
\usepackage{enumitem}
\usepackage{xcolor}
\usepackage{fancyhdr}
\usepackage{titlesec}
\usepackage{listings}
\usepackage{tikz}
\usetikzlibrary{positioning,arrows.meta,shapes.geometric,fit,backgrounds}
\usepackage[hidelinks,colorlinks=true,linkcolor=black!70,urlcolor=blue!60!black]{hyperref}

% Couleurs
\definecolor{primary}{HTML}{2563EB}
\definecolor{codebg}{HTML}{F3F4F6}

% En-têtes / pieds
\pagestyle{fancy}
\fancyhf{}
\lhead{\small DAT --- Chorus Path Finder (version limitée)}
\rhead{\small v1.0}
\cfoot{\thepage}
\renewcommand{\headrulewidth}{0.4pt}

% Style code
\lstset{
  basicstyle=\ttfamily\footnotesize,
  breaklines=true,
  frame=single,
  rulecolor=\color{black!20},
  backgroundcolor=\color{codebg},
  columns=fullflexible,
  keepspaces=true,
  showstringspaces=false,
  xleftmargin=4pt,
  framexleftmargin=4pt,
}

\titleformat{\section}{\Large\bfseries\color{primary}}{\thesection}{0.6em}{}
\titleformat{\subsection}{\large\bfseries}{\thesubsection}{0.6em}{}

\newcommand{\code}[1]{\texttt{\small #1}}

% ---------------------------------------------------------------------------
\begin{document}

% ===========================================================================
% Page de garde
% ===========================================================================
\begin{titlepage}
\centering
\vspace*{2.5cm}
{\Huge\bfseries\color{primary} Document d'Architecture Technique\par}
\vspace{0.6cm}
{\LARGE Chorus Path Finder --- \emph{Version limitée}\par}
\vspace{2cm}

\begin{tabular}{ll}
\toprule
\textbf{Projet} & Chorus Path Finder (gestion de factures Chorus Pro) \\
\textbf{Périmètre} & Version limitée (branche \code{feature/version-limitee}) \\
\textbf{Version du document} & 1.0 \\
\textbf{Date} & 10 juin 2026 \\
\textbf{Statut} & Diffusable \\
\textbf{Auteur} & Équipe Gigalis \\
\bottomrule
\end{tabular}

\vfill
{\small Document interne --- ne pas diffuser hors du périmètre projet.}
\end{titlepage}

% ===========================================================================
\tableofcontents
\newpage

% ===========================================================================
\section*{Historique des révisions}
\addcontentsline{toc}{section}{Historique des révisions}
\begin{longtable}{p{1.5cm} p{2.2cm} p{3cm} p{6.5cm}}
\toprule
\textbf{Version} & \textbf{Date} & \textbf{Auteur} & \textbf{Modifications} \\
\midrule
\endhead
1.0 & 10/06/2026 & Équipe Gigalis & Création initiale du DAT pour la version limitée. \\
\bottomrule
\end{longtable}

% ===========================================================================
\section{Introduction}

\subsection{Objet du document}
Ce document décrit l'architecture technique de la \textbf{version limitée} de
l'application \textbf{Chorus Path Finder}, une solution de gestion du cycle de
vie des factures émises et reçues via \textbf{Chorus Pro} (portail de
facturation électronique du secteur public français), exposé par la passerelle
\textbf{PISTE}. Il a pour but de fournir aux équipes techniques une vision
complète des composants, des technologies, des intégrations externes, du
modèle de données et des modalités de déploiement et d'exploitation.

\subsection{Périmètre}
La présente version est dite \emph{limitée} car elle restreint le périmètre
fonctionnel exposé par rapport à la version complète~:
\begin{itemize}[nosep]
  \item le socle « \textbf{Flux} » (dépôt de flux, PDF\,$\rightarrow$\,flux, dépôt G2B) est
        \textbf{masqué} dans l'interface (réactivable par configuration)~;
  \item l'extraction de factures PDF est \textbf{assistée par LLM} (vLLM)~;
  \item la page « Rejets / refus » propose une \textbf{analyse LLM du motif} avec
        suggestion d'action.
\end{itemize}
Le backend reste fonctionnellement complet~: seules des surfaces de l'interface
sont masquées.

\subsection{Acronymes et définitions}
\begin{longtable}{p{2.6cm} p{12.5cm}}
\toprule
\textbf{Terme} & \textbf{Définition} \\
\midrule
\endhead
DAT & Document d'Architecture Technique. \\
Chorus Pro & Portail public de facturation électronique (DGFiP). \\
PISTE & Plateforme d'Intermédiation des Services pour la Transformation de l'État (passerelle API + OAuth2). \\
DAT & Document d'Architecture Technique. \\
SIRET & Identifiant d'établissement (14 chiffres). \\
LLM & \emph{Large Language Model} (modèle de langage). \\
vLLM & Moteur d'inférence LLM exposant une API compatible OpenAI. \\
Docling & Service de conversion/structuration de documents (OCR / mise en page). \\
Avoir & Facture à montant négatif (note de crédit). \\
\bottomrule
\end{longtable}

% ===========================================================================
\section{Contexte et objectifs}

\subsection{Contexte}
L'application permet de~: synchroniser depuis Chorus Pro les fournisseurs,
clients publics et factures~; consulter et filtrer les factures émises et
reçues~; soumettre, recycler et compléter des factures~; et gérer les rejets.
La version limitée vise un \textbf{déploiement de démonstration / pré-production}
recentré sur le parcours « factures », tout en introduisant l'assistance LLM.

\subsection{Objectifs de la version limitée}
\begin{enumerate}[nosep]
  \item Réduire la surface fonctionnelle visible (masquage du socle Flux) sans
        supprimer de code, de façon \textbf{réversible}.
  \item Brancher un \textbf{LLM} (vLLM / \code{gpt-oss-20b}) sur l'extraction des
        factures PDF.
  \item Fournir, sur les factures rejetées / à recycler / suspendues, une
        \textbf{suggestion de traitement} générée par le LLM (avec repli
        heuristique).
\end{enumerate}

% ===========================================================================
\section{Vue d'ensemble de l'architecture}

L'application suit une architecture \textbf{3-tiers} conteneurisée~: un frontend
statique (SPA) servi par Nginx, une API applicative Python, et une base de
données PostgreSQL. Les intégrations externes (Chorus Pro via PISTE, vLLM,
Docling) sont appelées par l'API et, pour le LLM côté navigateur, directement
par le frontend.

\begin{center}
\begin{tikzpicture}[
  node distance=1.1cm and 1.4cm,
  box/.style={draw, rounded corners, align=center, minimum height=1cm, minimum width=2.6cm, fill=blue!5},
  ext/.style={draw, dashed, rounded corners, align=center, minimum height=0.9cm, minimum width=2.6cm, fill=orange!8},
  arr/.style={-{Latex}, thick}, every node/.style={font=\small}
]
\node[box] (nav) {Navigateur\\(SPA React)};
\node[box, right=of nav] (front) {Frontend\\Nginx + dist};
\node[box, right=of front] (api) {API\\FastAPI (uvicorn)};
\node[box, right=of api] (db) {PostgreSQL\\16};

\draw[-{Latex}, thick] (nav) -- node[above]{HTTP} (front);
\draw[-{Latex}, thick] (front) -- node[above]{\,/api\,} (api);
\draw[-{Latex}, thick] (api) -- node[above]{SQL} (db);

\node[ext, below=1.4cm of api] (piste) {Chorus Pro\\/ PISTE (OAuth2)};
\node[ext, left=of piste] (docling) {Docling\\(PDF)};
\node[ext, right=of piste] (vllm) {vLLM\\\code{gpt-oss-20b}};

\draw[-{Latex}, dashed] (api) -- (piste);
\draw[-{Latex}, dashed] (api) -- (docling);
\draw[-{Latex}, dashed] (nav.south) to[bend right=20] (vllm.north);
\end{tikzpicture}
\end{center}

\begin{itemize}[nosep]
  \item \textbf{Frontend} : SPA compilée (Vite) servie en statique par Nginx, qui
        proxifie \code{/api} vers l'API.
  \item \textbf{API} : FastAPI/uvicorn ; logique métier Chorus, synchronisation,
        extraction PDF, authentification.
  \item \textbf{Base} : PostgreSQL, schéma géré par Alembic.
  \item \textbf{LLM} : l'extraction appelle vLLM (API OpenAI) ; appel émis côté
        navigateur (proxy en développement).
\end{itemize}

% ===========================================================================
\section{Architecture applicative}

\subsection{Frontend (SPA)}
Application React mono-page, organisée par pages (\code{src/pages}), composants
réutilisables (\code{src/components}, dont la bibliothèque shadcn/ui) et une
couche d'accès données (\code{src/lib}, \code{src/data}). La récupération des
données passe par TanStack Query vers l'API \code{/api/db/\dots}. Le routage est
assuré par React Router.

Points structurants de la version limitée~:
\begin{itemize}[nosep]
  \item \code{src/lib/featureFlags.ts} : drapeau \code{FLUX\_ENABLED} (défaut
        \emph{off}) qui conditionne l'affichage du socle Flux.
  \item \code{src/lib/llm/} : configuration LLM (\code{config.ts}), client chat
        partagé (\code{client.ts}, \emph{provider-aware}), prompt d'extraction
        (\code{ollama.ts}).
  \item \code{src/lib/chorus/rejectionAnalysis.ts} et
        \code{src/components/chorus/RejectionSuggestionCard.tsx} : analyse et
        suggestion de traitement des rejets.
\end{itemize}

\subsection{Backend (API)}
Serveur FastAPI (\code{scripts/chorus\_submit\_server.py}) qui expose~:
\begin{itemize}[nosep]
  \item une \textbf{API de données} (\code{/api/db/\dots}) adossée à PostgreSQL
        (factures, fournisseurs, clients/services publics, synchronisation)~;
  \item un \textbf{pont Chorus / PISTE} (\code{/api/chorus/\dots},
        \code{/api/reception/\dots}) pour la soumission, le recyclage, la
        complétion, la consultation et le téléchargement~;
  \item des services transverses~: extraction PDF (\code{/api/extract-pdf-text}),
        jeton OAuth (\code{/api/chorus/oauth-token}), annuaire, SIRET.
\end{itemize}
La logique Chorus est centralisée dans \code{chorus\_python/} (client API,
synchronisation \code{db/sync.py}, ordonnanceur \code{db/scheduler.py},
configuration \code{config.py}, entités \code{db/entities.py}).

\subsection{Persistance}
PostgreSQL~16. Le schéma est versionné par Alembic (répertoire
\code{alembic/versions}). L'accès se fait via SQLAlchemy~2 et psycopg~3.

% ===========================================================================
\section{Stack technique}

\subsection{Frontend}
\begin{longtable}{p{4.5cm} p{4cm} p{6cm}}
\toprule
\textbf{Composant} & \textbf{Version} & \textbf{Rôle} \\
\midrule
\endhead
React / React DOM & 18.3 & Bibliothèque UI \\
Vite & 6.x & Bundler / serveur de dev \\
TypeScript & 5.8 & Typage statique \\
shadcn/ui + Radix UI & --- & Composants d'interface \\
Tailwind CSS & 3.4 & Styles utilitaires \\
TanStack Query & 5.x & Cache / requêtes serveur \\
React Router & 6.x & Routage SPA \\
react-hook-form + zod & --- & Formulaires + validation \\
recharts, sonner, lucide & --- & Graphes, toasts, icônes \\
\bottomrule
\end{longtable}

\subsection{Backend}
\begin{longtable}{p{4.5cm} p{4cm} p{6cm}}
\toprule
\textbf{Composant} & \textbf{Version} & \textbf{Rôle} \\
\midrule
\endhead
Python & 3.12 & Langage \\
FastAPI & $\geq$0.110 & Framework API \\
uvicorn & $\geq$0.29 & Serveur ASGI \\
SQLAlchemy & 2.x & ORM \\
Alembic & 1.13+ & Migrations de schéma \\
psycopg & 3.x & Pilote PostgreSQL \\
pydantic / pydantic-settings & 2.x & Validation / configuration \\
APScheduler & 3.10+ & Synchronisation planifiée \\
PyJWT + bcrypt & --- & Authentification (jetons + hash) \\
PyMuPDF, pdfplumber, pypdf & --- & Extraction PDF locale \\
requests / httpx & --- & Appels HTTP (PISTE, Docling, vLLM) \\
\bottomrule
\end{longtable}

\subsection{Infrastructure}
PostgreSQL~16 (alpine), Nginx~1.27 (alpine), Docker \& Docker Compose.

% ===========================================================================
\section{Modèle de données}

Le schéma principal (PostgreSQL) est composé des tables suivantes~:

\begin{longtable}{p{4cm} p{11cm}}
\toprule
\textbf{Table} & \textbf{Description} \\
\midrule
\endhead
\code{tenants} & Locataires (multi-tenant) ; tenant par défaut configurable. \\
\code{suppliers} & Fournisseurs synchronisés depuis Chorus. \\
\code{public\_customers} & Clients publics (structures destinataires). \\
\code{public\_services} & Services exécutants rattachés aux clients publics. \\
\code{invoices} & Factures (émission/réception) ; métadonnées Chorus, statuts, montants. \\
\code{invoice\_lines} & Lignes de facture. \\
\code{sync\_runs} & Journal des exécutions de synchronisation Chorus\,$\rightarrow$\,DB. \\
\code{chorus\_flux\_deposits} & Dépôts de flux et comptes rendus (module Flux). \\
\code{invoice\_events}, \code{api\_logs} & Module \emph{invoicepilot} (événements, journaux d'API). \\
\bottomrule
\end{longtable}

\paragraph{Point d'attention --- montants négatifs.}
La migration \code{20260609\_0006} \textbf{supprime les contraintes}
\code{total\_ht/total\_vat/total\_ttc $\geq$ 0} sur \code{invoices}. En effet,
Chorus Pro renvoie des \textbf{avoirs} et des factures \textbf{rejetées} à
montants négatifs~; ces contraintes bloquaient la synchronisation.

% ===========================================================================
\section{Intégrations externes}

\subsection{Chorus Pro via PISTE (OAuth2)}
L'API obtient un jeton OAuth2 (\code{oauth.piste.gouv.fr}) à partir d'un
\code{client\_id}/\code{client\_secret}, et émet les appels métier vers
\code{api.piste.gouv.fr} avec, en en-tête, le jeton \code{Bearer} et le compte
technique Chorus (\code{cpro-account}, base64 de \code{login:password}).
L'environnement (\textbf{qualif} / \textbf{prod}) est piloté par \code{CHORUS\_ENV} ;
les identifiants sont sélectionnés par suffixe (\code{CHORUS\_PROD\_*} /
\code{CHORUS\_QUALIF\_*}, avec repli sur \code{VITE\_CHORUS\_*}).

\subsection{vLLM (extraction LLM)}
Moteur d'inférence exposant une API compatible OpenAI
(\code{/v1/chat/completions}). Configuré via \code{GEN\_PROVIDER=vllm},
\code{GEN\_BASE\_URL}, \code{GEN\_MODEL} (\code{gpt-oss-20b}). Le client est
\emph{provider-aware} et gère également Ollama (\code{/api/chat}).

\subsection{Docling (conversion PDF)}
Service externe de structuration de documents PDF (\code{DOCLING\_SERVER\_URL}),
utilisé en complément de l'extraction locale (PyMuPDF / pdfplumber).

% ===========================================================================
\section{Flux fonctionnels clés}

\subsection{Synchronisation Chorus $\rightarrow$ DB}
Assurée par APScheduler (\code{db/scheduler.py})~: une \textbf{synchronisation
complète au démarrage} (\code{SYNC\_STARTUP\_FULL\_SYNC}, activée par défaut) et
une \textbf{synchronisation périodique} optionnelle
(\code{SYNC\_CRON\_ENABLED}, intervalle \code{SYNC\_CRON\_INTERVAL\_SECONDS}).
Le déclenchement manuel passe par \code{POST /api/db/sync/suppliers} puis
\code{/sync/invoices}. L'environnement utilisé est résolu via \code{CHORUS\_ENV}
ou l'\code{environment} transmis par le frontend.

\subsection{Extraction PDF assistée par LLM}
Le frontend extrait d'abord les champs par heuristiques
(\code{src/lib/chorus/extraction.ts})~; si un LLM est configuré
(\code{canUseLlmExtraction}), il \textbf{affine} le résultat via vLLM, avec
\textbf{repli automatique} sur les heuristiques en cas d'échec/timeout. Côté
backend, l'extraction repose sur PyMuPDF/pdfplumber et, optionnellement, Docling.

\subsection{Traitement des rejets (suggestion LLM)}
Sur les factures rejetées / à recycler / suspendues, le module
\code{rejectionAnalysis.ts} analyse le \textbf{motif Chorus} et propose une action
structurée~:
\begin{center}
\begin{tabular}{ll}
\toprule
\textbf{Action} & \textbf{Aiguillage} \\
\midrule
\code{RESOUMETTRE} & \code{/invoices/new?invoiceId=X\&correction=1} (pré-remplie) \\
\code{RECYCLER} & Fiche facture (formulaire de recyclage) \\
\code{COMPLETER} & Fiche facture (complétion d'une suspendue) \\
\code{CONTACTER\_CLIENT} & Information (blocage côté client) \\
\code{AUCUNE} & Information (ex.~doublon) \\
\bottomrule
\end{tabular}
\end{center}
La suggestion s'affiche dans le panneau de la liste « Rejets » et sur la fiche
facture. En l'absence de LLM, un \textbf{repli heuristique} (mots-clés) fournit
la même structure.

% ===========================================================================
\section{Spécificités de la version limitée}

\begin{longtable}{p{4cm} p{11cm}}
\toprule
\textbf{Aspect} & \textbf{Mise en œuvre} \\
\midrule
\endhead
Masquage du socle Flux & Drapeau \code{FLUX\_ENABLED} (\code{src/lib/featureFlags.ts}, défaut \emph{off}, surcharge \code{VITE\_FLUX\_ENABLED=true}). Masque le groupe de navigation « Flux » et redirige les routes flux. Backend inchangé. \\
Extraction LLM (vLLM) & Provider \code{vllm} ajouté (API OpenAI), variables \code{GEN\_*} exposées au build (\code{envPrefix}), appel branché en fin d'extraction. \\
Suggestion de rejets & Analyse LLM + repli heuristique ; bouton « Resoumettre » pré-rempli. \\
\bottomrule
\end{longtable}

% ===========================================================================
\section{Sécurité}

\begin{itemize}[nosep]
  \item \textbf{Authentification} : jetons JWT (PyJWT, secret \code{APP\_AUTH\_JWT\_SECRET},
        durée \code{APP\_AUTH\_TOKEN\_TTL\_SECONDS}). Comptes en mode fichier
        (\code{config/auth-users.json}, mots de passe \textbf{hachés bcrypt}) ou
        utilisateur unique par variables \code{APP\_AUTH\_EMAIL} /
        \code{APP\_AUTH\_PASSWORD\_HASH}.
  \item \textbf{Secrets} : \code{.env}, \code{.env.local} et
        \code{config/auth-users.json} sont \textbf{exclus du dépôt}
        (\code{.gitignore}) et fournis à chaque déploiement.
  \item \textbf{Chorus / PISTE} : identifiants OAuth2 et compte technique stockés
        en variables d'environnement, jamais committés.
  \item \textbf{CORS / réseau} : l'appel vLLM est émis côté navigateur ; en
        développement il transite par un proxy Vite (\code{/\_\_gen-proxy/vllm}).
\end{itemize}

% ===========================================================================
\section{Déploiement et exploitation}

\subsection{Topologie Docker}
\code{docker-compose.yml} (développement / déploiement simple) définit trois
services~:
\begin{longtable}{p{3cm} p{5.5cm} p{5.5cm}}
\toprule
\textbf{Service} & \textbf{Image / build} & \textbf{Port hôte (défaut)} \\
\midrule
\endhead
\code{postgres} & postgres:16-alpine & \code{POSTGRES\_HOST\_PORT} (5432) \\
\code{chorus-api} & \code{Dockerfile.backend} (uvicorn) & \code{CHORUS\_API\_HOST\_PORT} (8008) \\
\code{frontend} & \code{Dockerfile.frontend} (build Node $\rightarrow$ Nginx) & \code{APP\_HTTP\_PORT} (8081) \\
\bottomrule
\end{longtable}
\code{docker-compose.prod.yml} ajoute un service \code{migrate}
(\code{alembic upgrade head}) et un service de préparation des volumes.

\subsection{Lancement}
\begin{lstlisting}[language=bash]
# Construction + démarrage
docker compose up -d --build

# Migrations (base neuve)
docker compose run --rm chorus-api alembic upgrade heads

# Suivi des logs
docker compose logs -f chorus-api
\end{lstlisting}

\subsection{Configuration (variables principales)}
\begin{lstlisting}
# Environnement Chorus (backend) + badge (frontend, build)
CHORUS_ENV=prod
VITE_CHORUS_ENV=prod

# LLM (extraction)
GEN_PROVIDER=vllm
GEN_BASE_URL=http://<hote-vllm>:8003/v1
GEN_MODEL=openai/gpt-oss-20b

# Synchronisation
SYNC_CRON_ENABLED=true
SYNC_STARTUP_FULL_SYNC=true

# Ports hôtes (isolation multi-instances)
COMPOSE_PROJECT_NAME=chorus-limited
POSTGRES_HOST_PORT=5433
CHORUS_API_HOST_PORT=8009
APP_HTTP_PORT=8090
\end{lstlisting}

\subsection{Migrations}
Le graphe Alembic présente un \emph{branchpoint} (module \emph{invoicepilot})~;
sur une base neuve, utiliser \code{alembic upgrade heads} (au pluriel) pour
appliquer toutes les têtes.

% ===========================================================================
\section{Points d'attention et limites connues}

\begin{itemize}[nosep]
  \item \textbf{Compte technique PISTE partagé} : exécuter deux instances qui
        synchronisent simultanément avec le \emph{même} compte technique
        provoque des rejets d'authentification (HTTP 401) côté PISTE.
        Recommandation~: une instance de synchronisation active par compte
        technique~; les instances de test restent passives
        (\code{SYNC\_CRON\_ENABLED=false}, \code{SYNC\_STARTUP\_FULL\_SYNC=false}).
  \item \textbf{Badge environnement} : le libellé « Production / Sandbox » dépend
        de \code{VITE\_CHORUS\_ENV} \emph{figé au build} du frontend (lu dans le
        \code{.env} racine), indépendamment du backend.
  \item \textbf{Encodage des fichiers d'environnement} : les \code{.env} doivent
        être en \textbf{UTF-8} (un octet non-UTF-8 fait échouer le chargement
        \code{load\_dotenv} au démarrage).
  \item \textbf{Fichiers non versionnés} : \code{.env}, \code{.env.local} et
        \code{config/auth-users.json} doivent être fournis manuellement après un
        clone.
  \item \textbf{Appel LLM côté navigateur} : nécessite l'accessibilité réseau du
        serveur vLLM depuis le poste client (et la cohérence http/https pour
        éviter le \emph{mixed content}).
\end{itemize}

% ===========================================================================
\section{Annexes}

\subsection{Arborescence (extrait)}
\begin{lstlisting}
chorus-path-finder/
|- src/                      # Frontend React
|  |- pages/                 # Pages (InvoiceList, EmissionRejections, ...)
|  |- components/            # Composants (chorus/, layout/, ui/)
|  |- lib/
|  |  |- featureFlags.ts     # FLUX_ENABLED
|  |  |- llm/                # config.ts, client.ts, ollama.ts
|  |  |- chorus/             # extraction.ts, rejectionAnalysis.ts, ...
|- chorus_python/            # Backend (client Chorus, db, config)
|  |- db/                    # entities.py, sync.py, scheduler.py, api.py
|- scripts/chorus_submit_server.py  # Application FastAPI
|- alembic/versions/         # Migrations
|- docker-compose.yml / .prod.yml
|- Dockerfile.backend / Dockerfile.frontend
\end{lstlisting}

\subsection{Principaux points d'entrée API}
\begin{longtable}{p{6.5cm} p{8cm}}
\toprule
\textbf{Endpoint} & \textbf{Rôle} \\
\midrule
\endhead
\code{/api/db/sync/suppliers}, \code{/sync/invoices} & Synchronisation Chorus\,$\rightarrow$\,DB \\
\code{/api/db/invoices}, \code{/suppliers}, \dots & Lecture des données \\
\code{/api/chorus/submit}, \code{/recycle} & Soumission / recyclage \\
\code{/api/reception/consult}, \code{/process}, \code{/complete} & Réception \\
\code{/api/extract-pdf-text} & Extraction PDF \\
\code{/api/chorus/oauth-token} & Jeton OAuth PISTE \\
\bottomrule
\end{longtable}

\end{document}
